\documentclass[border=8pt]{standalone}
\usepackage{libertinus}
\usepackage{amsmath}
\usepackage{xcolor}
\usepackage{tikz}
\usetikzlibrary{positioning,arrows.meta,calc,fit,backgrounds,shapes.geometric}

\definecolor{ink}{HTML}{1E2227}
\definecolor{muted}{HTML}{5C6670}
\definecolor{cRGB}{HTML}{2F6DB5}
\definecolor{cIR}{HTML}{C0622B}
\definecolor{cFuse}{HTML}{3B8B66}
\definecolor{cTime}{HTML}{7450A8}
\definecolor{cNeutral}{HTML}{66707B}
\definecolor{cData}{HTML}{8C6A2F}
\definecolor{cAlert}{HTML}{A8384A}
\definecolor{cHead}{HTML}{B08A2E}
\color{ink}

\tikzset{
  blk/.style   = {rounded corners=3pt, align=center, font=\small, inner xsep=6pt, inner ysep=4.5pt,
                  minimum height=0.8cm, line width=0.5pt, text=ink,
                  execute at begin node={\hyphenpenalty=10000}},
  vis/.style   = {blk, fill=cRGB!14, draw=cRGB!65},
  thr/.style   = {blk, fill=cIR!15, draw=cIR!65},
  fus/.style   = {blk, fill=cFuse!16, draw=cFuse!70},
  scn/.style   = {blk, fill=cTime!14, draw=cTime!65},
  neu/.style   = {blk, fill=cNeutral!10, draw=cNeutral!55},
  hed/.style   = {blk, fill=cHead!16, draw=cHead!70},
  dat/.style   = {blk, fill=cData!12, draw=cData!60},
  io/.style    = {align=center, font=\small, text=ink},
  op/.style    = {circle, draw=cNeutral!70, fill=white, line width=0.5pt, inner sep=0pt,
                  minimum size=5.2mm, font=\small},
  lab/.style   = {font=\footnotesize\itshape, text=muted, align=center},
  note/.style  = {font=\footnotesize, text=muted, align=left},
  ar/.style    = {-{Stealth[length=2.1mm,width=1.6mm]}, line width=0.55pt, draw=ink!75},
  dar/.style   = {ar, dash pattern=on 2.4pt off 1.6pt},
  gar/.style   = {ar, draw=cData, dash pattern=on 2.4pt off 1.6pt},
  box/.style   = {rounded corners=6pt, fill=cNeutral!5, draw=cNeutral!45, line width=0.5pt},
  over/.style  = {preaction={draw=white, line width=3.2pt, -}},
}

\begin{document}
\begin{tikzpicture}[font=\small, M/.style={blk, minimum width=2.4cm, minimum height=0.66cm, inner ysep=1.5pt, font=\footnotesize}]
  \node[font=\small\scshape, text=muted, anchor=west] at (-4.6,11.0) {(a) one frame at one pyramid level};
  \node[io] (z) at (0,-0.6) {$Z_{\ell,t}$};
  \node[M, emb] (inp) at (0,0.6) {$1{\times}1$ conv $\to u_t$};
  \node[M, nrm] (dt) at (-1.6,2.0) {$\Delta_t = \mathrm{softplus}(\cdot)$};
  \node[M, nrm] (bc) at (1.6,2.0) {$B_t,\ C_t$};
  \node[op] (g1) at (-1.6,3.3) {$\times$};
  \node[io, anchor=east, font=\footnotesize] at (-2.3,3.3) {$\tfrac{\tau_t}{\tau_0}\cdot\rho_{\ell,t}$};
  \node[M, scn, minimum width=5.0cm, minimum height=0.9cm] (up) at (0,4.7)
     {$S_t = e^{\Delta_t A}\odot \mathcal{W}_{H_t}(S_{t-1}) + \Delta_t B_t u_t$};
  \draw[ar] (z) -- (inp);
  \draw[ar] (inp) -| (dt); \draw[ar] (inp) -| (bc);
  \draw[ar] (dt) -- (g1); \draw[ar] (g1) -- (g1 |- up.south);
  \draw[ar] (bc) -- (bc |- up.south);

  \node[M, nrm, minimum width=3.0cm] (wp) at (5.6,6.2) {warp $\mathcal{W}_{H_t}$, reset at cuts};
  \node[io] (sp) at (9.2,6.2) {$S_{t-1}$};
  \draw[ar] (sp) -- (wp); \draw[ar] (wp) |- ($(up.east)+(0,0.3)$);

  \node[M, neu, minimum width=3.6cm] (rd) at (0,6.2) {$y_t = C_t\cdot S_t + D\,u_t$};
  \node[M, neu, minimum width=3.6cm] (sp2) at (0,7.55) {DW $3{\times}3$, $1{\times}1$ conv (zero-init)};
  \node[op] (add) at (0,8.85) {$+$};
  \node[io] (zo) at (0,10.0) {$Z'_{\ell,t}$ to the neck};
  \draw[ar] (up) -- (rd); \draw[ar] (rd) -- (sp2); \draw[ar] (sp2) -- (add); \draw[ar] (add) -- (zo);
  \draw[ar] (z.west) -- ++(-3.8,0) |- (add.west);
  \draw[ar] ($(up.east)+(0,-0.3)$) -- ++(2.0,0) node[right, io] {$S_t$ to frame $t{+}1$};

  \node[note, text width=5.6cm, anchor=north west] at (3.6,2.6)
    {Each pixel keeps an $n$-dimensional state per channel. $\tau_t/\tau_0$ scales the step to the real time between frames;
     $\rho_{\ell,t}$, the higher of the two camera reliabilities from CMFM, lets an unreliable frame leave the memory
     untouched. The homography $H_t$ moves the memory with the camera, and a scene cut resets it.};

  \begin{scope}[xshift=13.2cm, yshift=4.6cm]
    \node[font=\small\scshape, text=muted, anchor=west] at (-0.6,6.4) {(b) unrolled over a stream};
    \foreach \k/\x/\lbl in {0/0/{t-2}, 1/3.2/{t-1}, 2/6.4/{t}} {
      \node[img] (fr\k) at (\x+0.9,0.3) {\includegraphics[width=1.7cm]{in_thermal.jpg}};
      \node[shp, below=3pt of fr\k] {frame $\lbl$};
      \node[M, fus, minimum width=2.4cm] (cf\k) at (\x+0.9,2.2) {CFFM-Net};
      \node[M, scn, minimum width=2.4cm] (ts\k) at (\x+0.9,3.7) {TSM};
      \node[M, hed, minimum width=2.4cm] (hd\k) at (\x+0.9,5.2) {neck, heads};
      \draw[ar] (fr\k) -- (cf\k); \draw[ar] (cf\k) -- (ts\k); \draw[ar] (ts\k) -- (hd\k);
    }
    \draw[ar] (ts0) -- node[lab, above=2pt] {$S$} (ts1);
    \draw[ar] (ts1) -- node[lab, above=2pt] {$S$} (ts2);
    \draw[ar] (ts2) -- ++(1.4,0) node[right, io] {$\cdots$};
    \node[note, text width=8.6cm, anchor=north west] at (-0.6,-1.5)
      {The fused features of each frame update a per-pixel memory, so a person who drops below the threshold in
       one frame is still remembered in the next. Phase 2 trains this on paired visible--thermal video.};
  \end{scope}
\end{tikzpicture}
\end{document}
